\section{Module \texttt{demo\_app.c}}

Le module \texttt{demo\_app.c} constitue le cœur de l'application de démonstration. Il orchestre l'assemblage des widgets via les structures de configuration du wrapper, gère l'état de l'interface et définit la logique de réaction aux événements utilisateur.

\subsection{Gestion de l'état : \texttt{showcase\_state}}

Pour maintenir une cohérence entre les différents composants (par exemple, mettre à jour un label de statut depuis un changement dans une entrée de texte), l'application utilise une structure d'état centralisée.

\begin{lstlisting}[language=C]
typedef struct {
  GtkWidget *status_label;
  GtkWidget *progress;
  GtkWidget *spinner;
  GtkWidget *switch_widget;
  bool spinner_running;
} showcase_state;
\end{lstlisting}

Cette structure est allouée dynamiquement lors du démarrage de l'application et passée en tant que \texttt{user\_data} à tous les callbacks.

\subsection{Callbacks et Logique Métier}

L'application implémente de nombreux callbacks pour démontrer l'interactivité. Par exemple, la mise à jour du statut lors de la saisie :

\begin{lstlisting}[language=C]
static void on_entry_changed(GtkEditable *editable, gpointer user_data) {
  showcase_state *state = user_data;
  char msg[96];
  g_snprintf(msg, sizeof(msg), "Input length: %zu",
             strlen(gtk_editable_get_text(editable)));
  set_status(state, msg);
}
\end{lstlisting}

Ou encore le pilotage asynchrone des boîtes de dialogue :

\begin{lstlisting}[language=C]
static void on_show_dialog(GtkButton *button, gpointer user_data) {
  showcase_state *state = user_data;
  GtkRoot *root = gtk_widget_get_root(state->status_label);
  alert_button_config buttons[] = {
      {.label = "Cancel", .appearance = "flat"},
      {.label = "Confirm", .appearance = "suggested"},
  };

  show_alert_dialog(&(alert_dialog_config){
      .parent = (root && GTK_IS_WINDOW(root)) ? GTK_WINDOW(root) : NULL,
      .message = "Proceed with wrapper-driven action?",
      .buttons = buttons,
      .button_count = G_N_ELEMENTS(buttons),
      .on_response = on_alert_result,
      .user_data = state,
  });
}
\end{lstlisting}

\subsection{Construction de l'Interface}

L'interface est découpée en "pages" intégrées dans un \texttt{GtkStack}. Chaque page est construite par une fonction dédiée (\texttt{make\_inputs\_page}, \texttt{make\_buttons\_page}, etc.) qui utilise les fonctions \texttt{create\_*} du wrapper.

La fonction \texttt{build\_showcase} assemble ces pages, ajoute une barre latérale de navigation (\texttt{GtkStackSidebar}), un menu principal (\texttt{GtkPopoverMenuBar}) et une barre de statut.

\subsection{APIs GTK et Logique Appliquée}

\begin{itemize}
    \item \textbf{\texttt{gtk\_application\_get\_active\_window()}} : Permet de récupérer la fenêtre actuellement focalisée pour servir de parent aux boîtes de dialogue surgissantes.
    \item \textbf{\texttt{gtk\_widget\_get\_root()}} : Utilisé pour remonter la hiérarchie des widgets à partir du label de statut afin de trouver la fenêtre racine.
    \item \textbf{\texttt{g\_snprintf()}} : Employé systématiquement pour formater les chaînes de caractères de statut de manière sécurisée (prévention des dépassements de tampon).
    \item \textbf{\texttt{g\_new0()}} : Alloue la structure \texttt{showcase\_state} en initialisant tous ses champs à zéro, garantissant qu'aucun pointeur n'est indéfini au démarrage.
    \item \textbf{\texttt{gtk\_window\_set\_child()}} : Définit le conteneur principal (\texttt{root}) comme contenu unique de la fenêtre d'application.
    \item \textbf{\texttt{g\_signal\_connect()}} : Point central de l'interactivité, reliant les signaux de bas niveau (clic, modification) aux fonctions de haut niveau de la démonstration.
\end{itemize}

\newpage
